\documentclass[a4paper,12pt]{article}
\usepackage[T2A]{fontenc}
\usepackage[utf8]{inputenc}
\usepackage[russian]{babel}
\usepackage{amsmath,amssymb,mathtools}
\usepackage{DejaVuSans}
\renewcommand{\familydefault}{\sfdefault}
\usepackage{icomma}
\usepackage[margin=22mm,top=24mm,bottom=23mm]{geometry}
\usepackage{microtype}
\usepackage{tikz}
\usetikzlibrary{arrows.meta,positioning,shapes.geometric,calc}
\usepackage{tabularx,booktabs}
\usepackage{enumitem}
\usepackage{hyperref}
\usepackage{parskip}
\usepackage{fancyhdr}
\usepackage{setspace}
\usepackage{xcolor}
\definecolor{accent}{HTML}{247678}
\definecolor{darkaccent}{HTML}{174C4D}
\definecolor{textgray}{HTML}{353B3B}
\definecolor{lightline}{HTML}{D8E2E1}
\hypersetup{hidelinks,pdftitle={Чек-лист: тригонометрические преобразования},pdfauthor={Лёвин Артём Александрович}}
\setlength{\parindent}{0pt}
\setlength{\headheight}{24pt}
\onehalfspacing
\setlist[itemize]{leftmargin=*,itemsep=1pt,topsep=3pt}
\setlist[enumerate]{leftmargin=*,itemsep=3pt,topsep=3pt}
\pagestyle{fancy}
\fancyhf{}
\fancyhead[L]{\footnotesize\color{darkaccent}Тригонометрические преобразования}
\fancyhead[R]{\footnotesize\color{textgray}23.10.2026}
\fancyfoot[C]{\footnotesize\color{textgray}\thepage}
\renewcommand{\headrulewidth}{0.3pt}
\renewcommand{\headrule}{\hbox to\headwidth{\color{lightline}\leaders\hrule height \headrulewidth\hfill}}
\newcommand{\checkitem}{\item[$\square$]}
\newcommand{\note}[1]{\par\smallskip{\color{darkaccent}\textbf{Обратите внимание.}} #1\par\smallskip}
\newcommand{\heading}[1]{\section*{\color{darkaccent}#1}}
\newcommand{\subheading}[1]{\subsection*{\color{accent}#1}}
\begin{document}
\thispagestyle{empty}
\begin{center}
\vspace*{26mm}
{\color{accent}\Large МАТЕМАТИКА \textbullet{} ПРОФИЛЬНЫЙ ЕГЭ}\par
\vspace{18mm}
{\color{darkaccent}\fontsize{31}{37}\selectfont\bfseries Чек-лист}\par
\vspace{9mm}
{\fontsize{20}{27}\selectfont Тригонометрические преобразования}\par
\vspace{6mm}
{\large От формул суммы до разложения уравнения на множители}\par
\vspace{18mm}
{\color{accent}\rule{75mm}{0.6pt}}\par
\vspace{13mm}
{\large\bfseries Николь Саркисьянц}\par
\vspace{7mm}
{\large 23 октября 2026 года}\par
\vfill
{\normalsize Лёвин Артём Александрович}\par
{\normalsize эксклюзивно для Николь Саркисьянц}\par
\vspace{25mm}
\end{center}
\begin{tikzpicture}[remember picture,overlay]
\draw[accent!60,line width=0.8pt]
([xshift=12mm,yshift=20mm]current page.south west)
.. controls ([xshift=45mm,yshift=28mm]current page.south west)
and ([xshift=-47mm,yshift=9mm]current page.south east)
.. ([xshift=-12mm,yshift=21mm]current page.south east);
\end{tikzpicture}
\clearpage

\heading{1. Что необходимо уметь}
Ключевая цель --- подобрать преобразование, после которого тригонометрическое уравнение раскладывается на множители или сводится к алгебраическому уравнению относительно одной тригонометрической функции.

\begin{itemize}
\checkitem Вычислять точные значения функций стандартных углов.
\checkitem Применять формулы синуса и косинуса суммы и разности.
\checkitem Узнавать сдвиг аргумента на $\frac{\pi}{2}$ и учитывать чётность функций.
\checkitem Использовать все три записи косинуса двойного угла.
\checkitem Преобразовывать произведение тригонометрических функций в сумму.
\checkitem Преобразовывать сумму или разность в произведение.
\checkitem Выбирать между стратегиями «один угол», «одна функция» и преобразованием произведения в сумму.
\checkitem Разлагать полученное уравнение на множители без потери корней.
\checkitem Получать полные серии решений простейших тригонометрических уравнений.
\end{itemize}

\subheading{Числовые углы и сдвиги}
\[
\begin{gathered}
\sin\frac{5\pi}{6}=\sin\left(\pi-\frac{\pi}{6}\right)=\frac12,
\qquad \sin\pi=0,\quad\cos\pi=-1,\\
\sin\frac\pi2=1,\qquad \cos\frac\pi2=0,\qquad
\cos\frac\pi6=\frac{\sqrt3}{2},\quad \sin\frac\pi6=\frac12.
\end{gathered}
\]
В частности, формула косинуса разности даёт
\[
\cos\left(2x-\frac\pi2\right)
=\cos2x\cos\frac\pi2+\sin2x\sin\frac\pi2
=\boxed{\sin2x}.
\]

\subheading{Формулы суммы и разности}
\[
\begin{aligned}
\sin(\alpha\pm\beta)&=\sin\alpha\cos\beta\pm\cos\alpha\sin\beta,\\
\cos(\alpha\pm\beta)&=\cos\alpha\cos\beta\mp\sin\alpha\sin\beta.
\end{aligned}
\]
\note{У синуса знак сохраняется; у косинуса перед произведением синусов знак меняется на противоположный.}

\subheading{Чётность и нечётность}
\[
\cos(-t)=\cos t,\qquad
\sin(-t)=-\sin t,\qquad
\tg(-t)=-\tg t,\qquad \ctg(-t)=-\ctg t.
\]
Последние две формулы применяются только в точках, где соответствующие функции определены.

\bigskip
\heading{2. Двойной угол: три формы косинуса}
\[
\boxed{\sin 2\alpha=2\sin\alpha\cos\alpha}
\]
\[
\begin{aligned}
\cos2\alpha
&=\cos^2\alpha-\sin^2\alpha,\\
&=1-2\sin^2\alpha,\\
&=2\cos^2\alpha-1.
\end{aligned}
\]
Для угла $4x$ вместо $\alpha$ подставляем $2x$:
\[
\sin4x=2\sin2x\cos2x,
\quad\cos4x=1-2\sin^22x=2\cos^22x-1.
\]
Из основного тригонометрического тождества следует
\[
\cos^22x=1-\sin^22x,\qquad
\sin^22x=1-\cos^22x.
\]
\note{Форму двойного угла выбирайте с учётом цели: если хотите оставить только синус, используйте $\cos4x=1-2\sin^22x$. Если только косинус, используйте $\cos4x=2\cos^22x-1$.}

\heading{3. Произведение превращается в сумму}
\[
\begin{aligned}
\sin\alpha\sin\beta
&=\frac{\cos(\alpha-\beta)-\cos(\alpha+\beta)}2,\\[2mm]
\cos\alpha\cos\beta
&=\frac{\cos(\alpha-\beta)+\cos(\alpha+\beta)}2,\\[2mm]
\sin\alpha\cos\beta
&=\frac{\sin(\alpha+\beta)+\sin(\alpha-\beta)}2.
\end{aligned}
\]
Порядок сомножителей можно менять, поскольку умножение коммутативно. В формуле $\sin\alpha\cos\beta$ важно верно определить, какой угол обозначен через $\alpha$.

\subheading{Пример}
\[
\begin{aligned}
\cos2x\sin4x
&=\sin4x\cos2x\\
&=\frac12\bigl(\sin(4x+2x)+\sin(4x-2x)\bigr)\\
&=\boxed{\frac12(\sin6x+\sin2x)}.
\end{aligned}
\]
\note{Если в произведении стоят разные углы, формулы перехода к сумме часто сокращают решение на несколько строк.}

\bigskip
\heading{4. Сумма превращается в произведение}
\[
\begin{aligned}
\sin\alpha+\sin\beta
&=2\sin\frac{\alpha+\beta}2\cos\frac{\alpha-\beta}2,\\[2mm]
\sin\alpha-\sin\beta
&=2\cos\frac{\alpha+\beta}2\sin\frac{\alpha-\beta}2,\\[2mm]
\cos\alpha+\cos\beta
&=2\cos\frac{\alpha+\beta}2\cos\frac{\alpha-\beta}2,\\[2mm]
\cos\alpha-\cos\beta
&=-2\sin\frac{\alpha+\beta}2\sin\frac{\alpha-\beta}2.
\end{aligned}
\]

\subheading{Что происходит с углами?}
В каждом множителе образуются \textbf{полусумма} или \textbf{полуразность} исходных углов. Деление на $2$ выполняется в обоих выражениях.
\[
\begin{aligned}
\sin6x-\sin2x
&=2\cos\frac{6x+2x}2\sin\frac{6x-2x}2\\
&=\boxed{2\cos4x\sin2x},\\[2mm]
\cos6x+\cos2x
&=2\cos\frac{6x+2x}2\cos\frac{6x-2x}2\\
&=\boxed{2\cos4x\cos2x}.
\end{aligned}
\]
\note{Типичная ошибка: вычислить $6x-2x=4x$, но забыть разделить полученную разность на $2$. Правильная полуразность равна $2x$.}

\clearpage
\heading{5. Один пример --- два полноценных решения}
Решите уравнение
\[
\boxed{\cos2x\sin4x-\sin\frac{5\pi}{6}\cos4x
=\cos\left(2x-\frac\pi2\right).}
\]
Сначала преобразуем числовой угол и правую часть:
\[
\cos2x\sin4x-\frac12\cos4x=\sin2x.
\tag{$*$}\label{eq:base}
\]

\subheading{Способ I. Один угол, одна функция}
Применяем $\sin4x=2\sin2x\cos2x$ и
$\cos4x=1-2\sin^22x$:
\[
2\sin2x\cos^22x-\frac12\bigl(1-2\sin^22x\bigr)=\sin2x.
\]
Теперь $\cos^22x=1-\sin^22x$:
\[
2\sin2x(1-\sin^22x)-\frac12+\sin^22x-\sin2x=0.
\]
Раскрываем скобки, приводим подобные и упорядочиваем по степеням:
\[
-2\sin^32x+\sin^22x+\sin2x-\frac12=0.
\]
Умножаем \emph{всё уравнение} на $-2$ и вводим $t=\sin2x$:
\[
4t^3-2t^2-2t+1=0,\qquad -1\le t\le1.
\]
Группируем и выносим общий множитель:
\[
\begin{aligned}
4t^3-2t^2-2t+1
&=2t^2(2t-1)-(2t-1)\\
&=(2t-1)(2t^2-1).
\end{aligned}
\]
Поэтому
\[
\sin2x=\frac12
\quad\text{или}\quad
\sin2x=\pm\frac{\sqrt2}{2}.
\]
Здесь $2\sin^22x-1=-\cos4x$, поэтому последние два значения равносильны условию $\cos4x=0$. Получаем две группы простейших уравнений:
\[
\boxed{\sin2x=\frac12\quad\text{или}\quad\cos4x=0.}
\]

\clearpage
\heading{6. Тот же пример: короткий путь}
\subheading{Способ II. Произведение $\rightarrow$ сумма $\rightarrow$ произведение}
Возвращаемся к уравнению~\eqref{eq:base}. Применяем формулу
\[
\sin4x\cos2x=\frac12(\sin6x+\sin2x).
\]
Следовательно,
\[
\frac12(\sin6x+\sin2x)-\frac12\cos4x=\sin2x.
\]
Умножаем на $2$, переносим слагаемые и приводим подобные:
\[
\sin6x-\sin2x-\cos4x=0.
\]
Преобразуем разность синусов:
\[
\sin6x-\sin2x=2\sin2x\cos4x.
\]
Итак,
\[
2\sin2x\cos4x-\cos4x=0
\iff \boxed{\cos4x(2\sin2x-1)=0}.
\]
Отсюда
\[
\cos4x=0\quad\text{или}\quad \sin2x=\frac12.
\]
Окончательно:
\[
\begin{aligned}
\cos4x=0&\iff 4x=\frac\pi2+\pi k
&&\iff x=\frac\pi8+\frac{\pi k}4,\\[2mm]
\sin2x=\frac12&\iff
\begin{cases}2x=\frac\pi6+2\pi n,\\2x=\frac{5\pi}6+2\pi n\end{cases}
&&\iff
\begin{cases}x=\frac\pi{12}+\pi n,\\x=\frac{5\pi}{12}+\pi n.\end{cases}
\end{aligned}
\]
Здесь $k,n\in\mathbb Z$.
\[
\boxed{x=\frac\pi8+\frac{\pi k}4\ (k\in\mathbb Z);\quad
x=\frac\pi{12}+\pi n\ \text{или}\ x=\frac{5\pi}{12}+\pi n\ (n\in\mathbb Z).}
\]
\note{Оба способа верны. Второй здесь короче: удачная формула превращает выражение в произведение почти сразу. Первый полезен, когда другие преобразования не приводят к факторизации.}

\clearpage
\heading{7. Как выбрать преобразование}
\vspace{2mm}
\centering
\begin{tikzpicture}[
  >=Latex,
  node distance=7mm and 9mm,
  start/.style={ellipse,draw=accent,very thick,align=center,text width=52mm,minimum height=9mm,fill=accent!6,font=\small},
  action/.style={rounded corners=2mm,draw=lightline,thick,align=center,text width=54mm,minimum height=12mm,font=\small,fill=white,inner sep=3mm},
  branch/.style={diamond,aspect=2.1,draw=accent,thick,align=center,text width=38mm,inner sep=1mm,font=\small,fill=accent!4},
  path/.style={-{Latex[length=2.5mm]},thick,color=darkaccent},
  every node/.style={text=textgray}
]
\node[start] (s) {Тригонометрическое уравнение};
\node[action,below=of s] (pre) {Числовые углы; формулы суммы/разности; сдвиги на $\pi/2$};
\node[branch,below=7mm of pre] (q) {Есть удобное произведение с разными углами?};
\node[action,below left=12mm and 2mm of q] (a) {Преобразовать произведение в сумму, затем сумму в произведение};
\node[action,below right=12mm and 2mm of q] (b) {Уменьшить углы; перейти к одной функции};
\coordinate (join) at ($(a.south)!0.5!(b.south)+(0,-7mm)$);
\node[action,below=6mm of join,text width=62mm] (f) {Упростить и разложить на множители};
\node[action,below=of f,text width=62mm] (z) {Приравнять множители к нулю; решить простейшие уравнения};
\node[start,below=of z,text width=55mm] (end) {Проверить ограничения; записать полные серии ответов};
\draw[path] (s)--(pre);
\draw[path] (pre)--(q);
\draw[path] (q.south) |- ([yshift=4mm]a.north) -- (a.north);
\draw[path] (q.south) |- ([yshift=4mm]b.north) -- (b.north);
\node[font=\scriptsize,anchor=east] at ([xshift=-2mm,yshift=8mm]a.north) {да};
\node[font=\scriptsize,anchor=west] at ([xshift=2mm,yshift=8mm]b.north) {нет};
\draw[path] (a.south) |- (join) -- (f.north);
\draw[path] (b.south) |- (join);
\draw[path] (f)--(z);
\draw[path] (z)--(end);
\end{tikzpicture}
\par
\clearpage

\heading{8. Типичные ошибки и самопроверка}
\begin{enumerate}[label=\textbf{\arabic*.}]
\item \textbf{Пропущено деление на два.}
При переходе $\sin6x-\sin2x$ к произведению используйте
$\frac{6x-2x}2=2x$, а не $4x$.
\item \textbf{Перепутан знак в формуле косинуса разности.}
$\cos(\alpha-\beta)=\cos\alpha\cos\beta+\sin\alpha\sin\beta$.
\item \textbf{Минус внутри функции обработан неверно.}
$\cos(-2x)=\cos2x$, но $\sin(-2x)=-\sin2x$.
\item \textbf{Вынесение общего множителя сделано частично.}
Из $2a^2b-a^2$ получается $a^2(2b-1)$: каждое слагаемое делится на вынесенный множитель.
\item \textbf{После преобразований нет одинаковых скобок.}
Не навязывайте факторизацию: вернитесь к выбору формулы и попробуйте другой путь.
\item \textbf{Уравнение умножено не целиком.}
При умножении на $-2$ изменяется знак каждого слагаемого. Например, $(-2)\cdot(-0{,}5)=1$.
\item \textbf{Неполный ответ.}
Уравнение $\sin2x=\frac12$ даёт две серии; уравнение $\cos4x=0$ --- одну серию.
\end{enumerate}

\subheading{Контрольный список перед сдачей решения}
\begin{itemize}
\checkitem Я заменил все числовые углы на точные значения.
\checkitem Я явно указал формулу, которую применяю.
\checkitem В полусуммах и полуразностях разделил на $2$ оба угла.
\checkitem Я проверил, что знак совпадает с формулой.
\checkitem При вынесении множителя я преобразовал все слагаемые.
\checkitem Я не делил на выражение, которое может равняться нулю.
\checkitem В замене $t=\sin\alpha$ или $t=\cos\alpha$ учтён отрезок $[-1;1]$.
\checkitem Все семейства решений записаны с целыми параметрами.
\end{itemize}

\clearpage
\heading{9. Тренировочный блок --- Путь А}
\textbf{10 заданий.} Выполняйте по порядку: от прямого применения формул к самостоятельному выбору метода. Записывайте ключевые преобразования. В заданиях с уравнениями требуется \textbf{полный ответ}, $k\in\mathbb Z$.

\begin{enumerate}[label=\textbf{\arabic*.},itemsep=7mm]
\item Вычислите:
$\displaystyle \sin\frac{5\pi}{6},\quad \cos\frac\pi2,\quad \sin\frac\pi2,\quad\cos\pi.$

\item Упростите выражение:
$\displaystyle \cos\left(2x-\frac\pi2\right)+\sin(-2x).$

\item Представьте с помощью двойного угла:
$\displaystyle 2\sin2x\cos2x$ и $\displaystyle 1-2\sin^22x$.

\item Представьте в виде суммы:
$\displaystyle \sin5x\cos x.$

\item Представьте в виде суммы:
$\displaystyle \cos4x\cos2x.$

\item Представьте в виде произведения:
$\displaystyle \sin7x-\sin3x.$

\item Представьте в виде произведения:
$\displaystyle \cos6x-\cos2x.$

\item Решите уравнение:
$\displaystyle \cos6x+\cos2x=0.$

\item Разложите на множители:
$\displaystyle 4\sin^32x-2\sin^22x-2\sin2x+1.$

\item Решите двумя способами --- через двойной угол и через переход от произведения к сумме:
\[
\cos2x\sin4x-\sin\frac{5\pi}{6}\cos4x
=\cos\left(2x-\frac\pi2\right).
\]
\end{enumerate}

\clearpage
\heading{10. Ответы для самопроверки}
Сначала доведите решение самостоятельно до конца; затем сравните ответ.

\renewcommand{\arraystretch}{1.6}
\begin{tabularx}{\linewidth}{@{}>{\bfseries}p{11mm}X@{}}
\toprule
№ & Ответ \\
\midrule
1 & $\frac12;\quad 0;\quad 1;\quad -1$. \\
2 & $0$. \\
3 & $\sin4x$;\quad $\cos4x$. \\
4 & $\frac12(\sin6x+\sin4x)$. \\
5 & $\frac12(\cos6x+\cos2x)$. \\
6 & $2\cos5x\sin2x$. \\
7 & $-2\sin4x\sin2x$. \\
8 & $x=\frac\pi8+\frac{\pi k}{4}$ или $x=\frac\pi4+\frac{\pi k}{2}$, $k\in\mathbb Z$. \\
9 & $(2\sin2x-1)(2\sin^22x-1)$. \\
10 & $x=\frac\pi8+\frac{\pi k}{4}$; $x=\frac\pi{12}+\pi k$; $x=\frac{5\pi}{12}+\pi k$, $k\in\mathbb Z$. \\
\bottomrule
\end{tabularx}

\vspace{12mm}
\subheading{Смысл урока в трёх тезисах}
\begin{itemize}
\item Формулы выбираются по \textbf{структуре выражения}: сумма, разность, произведение, кратные углы.
\item Две полезные стратегии --- \textbf{один угол и одна функция} либо \textbf{произведение $\rightarrow$ сумма $\rightarrow$ произведение}.
\item Удобное преобразование узнаётся по \textbf{появлению общего множителя}: после этого уравнение часто решается за несколько строк.
\end{itemize}
\vfill
{\centering\color{textgray}\footnotesize Лёвин Артём Александрович эксклюзивно для Николь Саркисьянц\par}
\end{document}
